% =====================================================================
%  mcenerney-lens-handout.tex
%  Handout 1 -- Expert Academic Writing Methods
%
%  Source: L. McEnerney, "Leadership Lab: The Craft of Writing
%          Effectively", University of Chicago Writing Program,
%          recording posted 26 June 2014.
%  Pairs with Handout 2 (effective-writing): this handout checks a
%  paper's readers, problem and positioning; Handout 2 checks a single
%  paragraph against the same reader-value idea.
%
%  REPO NOTES
%  - The wording of the first four philosophy points, M1-M5, the
%    pitfalls and the Bonus Tip is the original handout's. The action
%    items (L1-L12) were rewritten on 6 Nov 2026 (see below); the
%    original questions are kept inside them.
%  - Added from a review of the lecture transcript (25 Oct 2026):
%    the fifth philosophy point, M6-M10 and L9-L13. They are tagged
%    ADDED in the PDF; set \markaddedfalse below to hide the tags.
%  - Also added: Sections 5-8 (using the lens, worked example,
%    worksheet, report format). Their code sits between the
%    ADDITIONS markers, in the preamble and before the References.
%  - Method IDs (M1 ... M10) and lens-field IDs (L1 ... L12) are stable
%    from the 6 Nov 2026 revision on.
%    The mcenerney-lens skill and its reports cite them, so never
%    renumber: append new IDs at the end instead. To print plain
%    numbers instead of M1/L1, edit \methodid and \lensid below.
%  - Time stamps (h:mm:ss) point to the recorded lecture so every
%    point can be traced back to where it is taught.
%  - Compiles with pdfLaTeX on Overleaf (TeX Live 2020 or later).
%    For A4 paper, change "letterpaper" in the geometry line.
%  - 25 Oct 2026, settled with Handout 2 for the skills repo: locations
%    are written P2 S3 (no pilcrow); L12 has a status cutoff; R5 caps
%    quotes; the L2 word list matches the compiled PDF (no "important");
%    the report format matches the mcenerney-lens skill; a "With
%    Handout 2" note after the Bonus Tip says what each handout leaves
%    to the other. In the worked example, P1 S3's "it drains" became
%    "the closure drains", because "it" first read as backscatter.
%  - 6 Nov 2026, action items restructured before the repo goes public
%    (the one-time renumbering). The lens reads the whole paper (R6);
%    readers and communities merged (L1); value words quoted in their
%    sentence, with reusable frames (L2); new fields for value offered
%    (L3), doubts answered (L7) and pitfalls (L10); fields L1-L10 cover
%    the whole paper; the problem in three parts (L4); "move" and
%    "positioning" renamed (L5 stance, L8 contribution), stance placed
%    right after the problem;
%    explanation order folded into structure (L9); figures became
%    locations; each field names the methods it draws on; R2 locations
%    use the sentence's first words; R5 allows one sentence per finding.
%    Old -> new: L1+L9 -> L1; L2+L8 -> L2; L3 -> L4; L4 -> L6; L5 -> L8;
%    L6 -> L5; L7+L10 -> L9; L11 -> figures as locations; L12 -> L11;
%    L13 -> L12.
%  - 6 Nov 2026, later: reports label fields by number and name
%    ("1- Readers"), not by ID; L10 became four named checks, each
%    tied to a Common Pitfall; report cells allow about 60 words.
% =====================================================================

\documentclass[11pt]{article}

% ---------- Fonts and encoding ---------------------------------------
\usepackage[T1]{fontenc}
\usepackage[utf8]{inputenc}
\usepackage{charter}                 % body: Bitstream Charter
\usepackage[scaled=0.90]{helvet}     % sans: Helvetica clone for headings
\usepackage{microtype}

% ---------- Page layout ----------------------------------------------
\usepackage[letterpaper,margin=0.9in,headheight=14pt]{geometry}
\usepackage{xcolor}
\usepackage{enumitem}
\usepackage{titlesec}
\usepackage{fancyhdr}
\usepackage[most]{tcolorbox}
\usepackage[hidelinks]{hyperref}
\hypersetup{pdftitle={Expert Academic Writing Methods (Handout 1)}, pdfauthor={Mohammed Meligy}}

\urlstyle{same}
\setlength{\parindent}{0pt}
\setlength{\parskip}{5pt}

% ---------- Colours (shared with Handout 2) --------------------------
\definecolor{ink}{HTML}{1F2328}
\definecolor{accent}{HTML}{1F4E79}
\definecolor{muted}{HTML}{5B6B7D}
\definecolor{rulec}{HTML}{C9D3DF}
\definecolor{fill}{HTML}{F2F5F9}
\color{ink}

% ---------- Headings -------------------------------------------------
\titleformat{\section}{\sffamily\large\bfseries\color{accent}}{\thesection}{0.6em}{}
\titleformat{\subsection}{\sffamily\normalsize\bfseries\color{ink}}{}{0em}{}
\titlespacing*{\section}{0pt}{16pt}{6pt}
\titlespacing*{\subsection}{0pt}{10pt}{3pt}

% ---------- Header and footer ----------------------------------------
\pagestyle{fancy}
\fancyhf{}
\renewcommand{\headrulewidth}{0.4pt}
\renewcommand{\headrule}{\color{rulec}\hrule height\headrulewidth}
\fancyhead[L]{\sffamily\footnotesize\color{muted}Expert Academic Writing Methods}
\fancyhead[R]{\sffamily\footnotesize\color{muted}mcenerney-lens}
\fancyfoot[C]{\sffamily\footnotesize\color{muted}\thepage}

% ---------- Small helpers --------------------------------------------
\newcommand{\tsin}[1]{\mbox{\sffamily\footnotesize\color{muted}[#1]}} % time stamp inside a sentence
% Time stamp ending an item: sits flush right on the last line, or on its own line if it does not fit.
\newcommand{\ts}[1]{{\unskip\nobreak\hfil\penalty50\hskip0.8em\hbox{}\nobreak\hfil\tsin{#1}\parfillskip=0pt\par}}
\newcommand{\methodid}[1]{M#1}                                    % method ID shown on cards
\newcommand{\lensid}[1]{L#1}                                      % lens-field ID shown on action items
\newif\ifmarkadded \markaddedtrue   % \markaddedfalse hides the ADDED tags
\newcommand{\added}{\ifmarkadded\,{\setlength{\fboxsep}{1.2pt}\fcolorbox{rulec}{white}{\sffamily\tiny\bfseries\color{accent}ADDED}}\fi}

% >>> ADDITIONS-PREAMBLE BEGIN (worked example, worksheet, report format)
\usepackage{array,tabularx,booktabs}
\usepackage{needspace}
\definecolor{strongc}{HTML}{2E6B3F}
\definecolor{partialc}{HTML}{9A6700}
\definecolor{missingc}{HTML}{A33A2B}
\newcommand{\badge}[2]{{\setlength{\fboxsep}{1.6pt}\colorbox{#1}{\sffamily\scriptsize\bfseries\color{white}\strut#2}}}
\newcommand{\stS}{\badge{strongc}{STRONG}}
\newcommand{\stP}{\badge{partialc}{PARTIAL}}
\newcommand{\stM}{\badge{missingc}{MISSING}}
\newcommand{\stN}{\badge{muted}{N/A}}
\newcommand{\sn}[1]{\textsuperscript{\sffamily\bfseries\color{accent}S#1}\,}  % sentence label in example text
\newcommand{\pn}[1]{{\sffamily\bfseries\color{accent}P#1}\enspace}           % paragraph label in example text
\newcommand{\loc}[1]{\leavevmode{\sffamily\footnotesize\color{accent}#1}}       % location in tables
% <<< ADDITIONS-PREAMBLE END

% Lists
\setlist{itemsep=3pt,parsep=0pt,topsep=3pt}
\setlist[itemize]{leftmargin=1.4em}
\newlist{lensfields}{description}{1}
\setlist[lensfields]{style=multiline,leftmargin=2.6em,font=\sffamily\bfseries\color{accent},itemsep=3pt}

% Method card:  \begin{method}{ID}{Title}{time stamps} ... \end{method}
\newtcolorbox{method}[3]{%
  enhanced,
  colback=white, colframe=rulec, boxrule=0.6pt, arc=2pt,
  colbacktitle=fill, coltitle=ink, toptitle=3pt, bottomtitle=3pt,
  left=8pt, right=8pt, top=5pt, bottom=5pt,
  fonttitle=\sffamily\bfseries,
  title={\textcolor{accent}{#1}\quad #2\hfill{\mdseries\footnotesize\color{muted}#3}},
  before skip=9pt, after skip=9pt}

% Plain shaded panel for notes
\newtcolorbox{panel}[1][]{%
  enhanced, breakable, colback=fill, colframe=fill, arc=2pt,
  left=10pt, right=10pt, top=7pt, bottom=7pt,
  before skip=8pt, after skip=8pt, #1}

% =====================================================================
\begin{document}

% ---------- Title block ----------------------------------------------
{\sffamily\footnotesize\bfseries\color{muted}HANDOUT 1 \textperiodcentered{} MCENERNEY-LENS}\par\vspace{2pt}
{\sffamily\LARGE\bfseries Expert Academic Writing Methods}\par\vspace{3pt}
{\sffamily\small\bfseries Mohammed Meligy}{\sffamily\small\color{muted}\ \textperiodcentered{} Mechanical and Aerospace Engineering, Arizona State University \textperiodcentered{} \textcopyright{} 2026, CC BY 4.0}\par\vspace{3pt}
{\small\color{muted}Compiled from \emph{The Craft of Writing Effectively}, a Leadership Lab lecture by Larry McEnerney
(University of Chicago Writing Program); recording posted 26 June 2014. Unofficial study handout: time stamps
such as \tsin{31:34} point to the recording.}\par
\vspace{4pt}{\color{rulec}\rule{\linewidth}{0.6pt}}

% =====================================================================
\section{Writing Philosophy Overview}

\begin{itemize}
  \item Academic writing is not about rules; it's about \textbf{readers}. \ts{2:59--3:54}
  \item The purpose of writing is not to convey what you think \emph{but to change what your readers think}.
        \ts{21:22--22:07}
  \item Your writing succeeds only if it is perceived as \textbf{valuable} by a specific community of readers.
        \ts{14:08--16:24}
  \item Value is defined \emph{not} in the content itself but by how readers use or respond to it.
        \ts{15:27--16:31}
  \item Knowledge is a conversation that moves through time, not a pile that grows with every new fact. The
        community decides what counts as knowledge, and ideas also get left behind; new information it does not
        care about is not knowledge.\added \ts{28:03--30:51}
\end{itemize}

% =====================================================================
\section{Core Methods for Effective Academic Writing}

\begin{method}{\methodid{1}}{Value-Oriented Writing}{22:13--34:10}
\begin{itemize}
  \item Ask not \emph{what do I want to say?} but \emph{why should my readers care?} \ts{22:13--22:38}
  \item Use language that signals value: words like \emph{nonetheless}, \emph{however}, \emph{although},
        \emph{inconsistent}, \emph{anomaly}, \emph{reported}, \emph{widely}, \emph{accepted}. \ts{31:28--32:21}
  \item Spend 15 minutes a week reading articles in your field and \textbf{circling value-signaling words}.
        Build your own word bank. \ts{32:27--34:10}
\end{itemize}
\end{method}

\begin{method}{\methodid{2}}{Know Your Readers}{16:37--17:10; 34:46--35:20; 1:00:57}
\begin{itemize}
  \item Understand the beliefs, assumptions, and doubts of your specific disciplinary audience. \ts{34:46--35:20}
  \item Avoid generic writing aimed at abstract or standardized readers (e.g., SAT/AP). \ts{16:45--17:10}
  \item Learn how your target journals express \textbf{costs} (what's at stake) and \textbf{benefits} (what is
        gained). \ts{1:00:57--1:02:37}
\end{itemize}
\end{method}

\begin{method}{\methodid{3}}{Constructing Problems Effectively}{37:24--38:40; 55:19--59:05}
\begin{itemize}
  \item Shift from background or definitions to \textbf{instability} or \textbf{tension}. \ts{55:19--56:25; 58:42}
  \item Build problems that \textbf{matter to your readers}, not problems of personal interest. \ts{56:45--58:07}
  \item Problems should challenge assumptions or practices in the field. \ts{37:24--38:40}
\end{itemize}
\end{method}

\begin{method}{\methodid{4}}{Argument vs.\ Explanation}{20:18--24:48; 40:23--40:48}
\begin{itemize}
  \item Do not \emph{explain to show what you know}; \textbf{argue to change what your readers think}.
        \ts{20:18--20:57; 40:23}
  \item Explanation is inward-facing. Argumentation is outward-facing. \ts{20:57--21:17}
  \item Predict and overcome reader doubts. Your goal is not self-revelation but persuasion.
        \ts{35:09; 40:29--40:48}
\end{itemize}
\end{method}

\begin{method}{\methodid{5}}{Literature Reviews as Problem Builders}{1:03:14--1:10:45}
\begin{itemize}
  \item A professional literature review enriches the \textbf{problem}, rather than listing chronologically who
        said what. \ts{1:04:52--1:05:23}
  \item Highlight tensions, contradictions, or conflicts between works. \ts{1:05:23--1:06:43}
  \item Avoid the \emph{gap model} (``this hasn't been done''). Prefer the \textbf{error model} (``this has been
        done, but there's a problem''). \ts{1:08:47--1:10:45}
\end{itemize}
\end{method}

\begin{method}{\methodid{6}}{Two Kinds of Value Words\added}{35:27--37:12; 53:35--54:52}
\begin{itemize}
  \item \textbf{Community words} such as \emph{widely}, \emph{accepted} and \emph{reported} show that a community
        of readers already cares about the question. \ts{35:27--35:55}
  \item \textbf{Instability words} such as \emph{nonetheless}, \emph{however}, \emph{although}, \emph{inconsistent}
        and \emph{anomaly} show that something in that community's understanding is unstable. They create tension,
        whatever their part of speech. \ts{53:35--54:52}
  \item \textbf{Flow words} such as \emph{because}, \emph{and}, \emph{if} and \emph{unless} make prose flow but
        create no value. Don't count them as value words. \ts{36:15--37:12}
  \item \emph{Rarely} (``has rarely been scrutinized'') is often code for a gap claim. \ts{1:10:53--1:10:59}
\end{itemize}
\end{method}

\begin{method}{\methodid{7}}{Value Before Explanation\added}{19:52--25:13}
\begin{itemize}
  \item Your readers' professional function is to challenge what you say: nothing is accepted as knowledge until
        someone competent has challenged it. \ts{23:16--24:27}
  \item Explanation belongs only after value has been established and persuasion has begun. \ts{24:27--24:48}
  \item Don't open with a definition. It shows a teacher that you know it; it gives a professional reader
        nothing. \ts{25:00--25:13}
  \item When readers say your work doesn't seem important, or that they didn't understand it, don't respond by
        explaining more. \ts{19:52--20:23}
\end{itemize}
\end{method}

\begin{method}{\methodid{8}}{Openings That Show a Problem\added}{55:19--1:00:37; 1:07:55}
\begin{itemize}
  \item Avoid the ``martini glass'' opening: a generalization, then a thesis, then a generalization again.
        \ts{55:47--56:19}
  \item Readers scan an opening for instability. Language of continuity and consistency there interferes with
        their reading, and they can stop two paragraphs in. \ts{59:05--1:00:37}
  \item A thesis is a solution only if readers can see the problem it solves, so the problem must be stated.
        \ts{58:12--58:36}
  \item Locate the problem in something readers care about: for academic readers, something they want to
        understand; for other readers, something they want to fix. \ts{57:22--58:07}
  \item When readers ask for ``more background,'' they usually mean they can't see why it matters. The background
        they need is a problem. \ts{1:07:55--1:08:21}
\end{itemize}
\end{method}

\begin{method}{\methodid{9}}{Saying ``You're Wrong'' in the Community's Code\added}{39:05--44:13}
\begin{itemize}
  \item ``I have something to add'' invites the reply ``we don't care.'' Showing that the community has
        something wrong gives it a reason to listen. \ts{41:16--41:51}
  \item Every community has a polite code for ``you're wrong,'' often deferential: you have advanced this field,
        but there is one thing here that is wrong. Learn how your journals say it. \ts{39:05--40:14; 43:50--44:02}
  \item Break the code and you get slapped down; skip the move altogether and you get rejected.
        \ts{44:02--44:13}
\end{itemize}
\end{method}

\begin{method}{\methodid{10}}{Raising the Costs in a Literature Review\added}{1:04:21--1:06:43; 1:09:49}
\begin{itemize}
  \item Credibility and acknowledging the community's work are real functions of a literature review too.
        \ts{1:04:21--1:04:41}
  \item The strongest reviews raise the costs: the tension matters because the community has kept building on
        it, so ``we'' are wrong too, not only earlier authors. \ts{1:06:05--1:06:43}
  \item Gap claims assume knowledge is bounded, like a crossword puzzle. If knowledge is open-ended, filling one
        gap leaves countless others, so the claim adds little. \ts{1:09:49--1:10:45}
\end{itemize}
\end{method}

% =====================================================================
\section{Common Pitfalls to Avoid}

\begin{itemize}
  \item Thinking writing is about expressing your thoughts---\allowbreak it's about changing others' thoughts.
        \ts{21:30--22:07; 49:22}
  \item Starting with personal curiosity or background instead of the reader's problem. \ts{56:45--57:22; 1:08:07}
  \item Believing that ``originality'' guarantees value. \textbf{Value comes from relevance to the reader.}
        \ts{25:26--30:58}
  \item Assuming your writing will ``preserve ideas forever.'' The goal is participation, not permanence.
        \ts{46:45--48:58}
\end{itemize}

% =====================================================================
\section{Action Items: Reading an Article in Your Field}

Read the whole article, not only the introduction: problems, stakes, value and doubts are often stated again, or
stated more sharply, in later sections (R6). Fields L1--L10 cover the whole paper; for each, record what the text
does and where, in whichever section it appears. \emph{Draws on} names the methods behind each field.

\begin{lensfields}
  \item[\lensid{1}] \textbf{Identify the readers:} Two parts. \emph{Research community}: whose understanding does
        the paper set out to change? Interdisciplinary work has several communities, and each defines problems and
        judges arguments differently. \emph{Wider readers}: who could use the work to fix something, such as
        practitioners or engineers? Note the journal it is published in. Draws on M2, M8.
        \ts{34:46; 57:22--58:07; 1:12:19--1:14:04}
  \item[\lensid{2}] \textbf{Circle value words, in context:} Quote each value word in its sentence, with the word in
        bold. Sort it as an instability word (\emph{however}, \emph{nonetheless}, \emph{inconsistent},
        \emph{problematic}, \emph{anomaly}) or a community word (\emph{widely}, \emph{accepted}, \emph{reported}).
        Flow words don't count; \emph{gap} and \emph{rarely} often signal a gap claim. Where a sentence is worth
        reusing, add its frame: the sentence with the paper's specifics replaced by slots. Draws on M1, M6 and the
        Bonus Tip. \ts{31:34--34:10; 35:27--37:12; 53:35--54:52}
  \item[\lensid{3}] \textbf{Name the value offered:}\added{} What do the authors say readers gain, overall? Quote the
        sentence that says it. A type label, such as verification, method, benchmark, correction, theory or
        practical use, is your interpretation, so mark it as one. Draws on M1, M2.
        \ts{15:27--16:31; 1:00:57--1:02:37}
  \item[\lensid{4}] \textbf{Find the problem:} State it in three parts, each with its location. \emph{Current
        state}: what the community does or believes. \emph{Where it breaks}: each approach or assumption, and how it
        fails. \emph{What remains open}: the question the paper answers. Is the problem explicit or implicit?
        Draws on M3, M8. \ts{56:25--59:05}
  \item[\lensid{5}] \textbf{Read the stance toward prior work:} Who do the authors say got the problem wrong: the
        dominant view, rival approaches, or the community itself? Do they say it directly, or deferentially,
        crediting the work before naming its limit? Draws on M9, M10. \ts{39:05--44:13; 1:16:43}
  \item[\lensid{6}] \textbf{Find the stakes:} What does the text say it costs to ignore the problem, and what is
        gained by solving it? Note whether the cost is stated as a cost to these readers (their predictions,
        results or decisions) or only as a fact about the world. Draws on M2, M10.
        \ts{1:00:57--1:02:21; 1:06:05--1:06:43}
  \item[\lensid{7}] \textbf{Find the doubts answered:}\added{} Which reader doubts does the paper anticipate, such as
        limitations, validation or ``one might argue,'' and where does it answer them? Draws on M4, M7.
        \ts{23:16--24:27; 35:09; 40:29--40:48}
  \item[\lensid{8}] \textbf{Find the contribution:} What does the paper offer in response to the problem? Is it named
        as something readers can use (a tool, a test, a method) or only as an activity (``X is analyzed'')?
        Draws on M1, M8. \ts{58:12--58:36}
  \item[\lensid{9}] \textbf{Map the structure:} Does the paper build from instability: the problem, then prior work
        that sharpens it, then the contribution? Do the discussion and conclusion return to the problem? Does
        explanation (definitions, background, method) come only
        after value is established? Does the literature review build the problem or list who said what, and does it
        argue by the error model or the gap model? Draws on M3, M5, M7, M8.
        \ts{24:27--25:13; 55:19--56:19; 1:04:52--1:10:45}
  \item[\lensid{10}] \textbf{Check the pitfalls:}\added{} Run four checks, each drawn from a Common Pitfall
        (Section~3, in brackets), and report each as found, quoted with its location, or not found.
        \begin{itemize}[nosep,leftmargin=1.2em]
          \item \emph{Background opening} (starting with background): the introduction's first paragraph has no
                instability word. Either way, note where the first one appears.
          \item \emph{The writers' interest as the reason} (starting with personal curiosity): a sentence gives the
                authors' interest or aim (``we are interested in,'' ``this study aims to explore'') before the
                introduction states any problem.
          \item \emph{Originality as the value} (believing originality guarantees value): ``novel,'' ``new,''
                ``first'' or ``for the first time'' is offered as the reason to read, and the same paragraph does not
                say what it changes for readers.
          \item \emph{Results without consequence} (expressing your thoughts instead of changing the readers'): a
                results, discussion or conclusion paragraph says what a figure, table or case shows, and neither it
                nor the next paragraph says what that changes for the problem (L4) or for the readers.
        \end{itemize}
        The fourth pitfall, writing to preserve ideas, is a belief that doesn't show in a text, so it has no
        check. Draws on M4, M8. \ts{21:30--22:07; 25:26--30:58; 56:45--57:22}
\end{lensfields}

\subsection{For your own draft (self-check)}

\begin{lensfields}
  \item[\lensid{11}] \textbf{Count value words:}\added{} Can you circle about ten value words in your first two
        paragraphs? If not, go back to your word bank. For a self-check status, count distinct community and
        instability words (M6; flow words don't count): 8 or more is STRONG, 4--7 PARTIAL, 3 or fewer MISSING.
        Draws on M1, M6. \ts{34:04--34:18}
  \item[\lensid{12}] \textbf{Revise the opening last:}\added{} Does your first sentence state the problem the paper
        ended up solving? Check it against L4: if a later section states the problem more sharply, the opening
        should carry it. Writers often find their problem late, sometimes only from a final reader. Draws on M8.
        \ts{1:18:46--1:19:39}
\end{lensfields}

\textbf{Figures and tables} are locations, not a separate field. When a figure shows the problem, the stakes or the
value, for example the accepted view failing, cite it in that field. \tsin{1:14:53--1:15:20}

\begin{panel}
\textbf{Bonus Tip:} Maintain a personal glossary of high-value vocabulary and sentence structures from journal
articles. \ts{33:45--34:10}
\end{panel}

\textbf{With Handout 2.} On one draft, this handout owns the opening's problem, stance, stakes and contribution (L4,
L5, L6, L8, L12); Handout 2 (\emph{effective-writing}) then checks only the opening's sentences (EW-5--EW-10), so the
opening is not flagged twice.

% >>> ADDITIONS-BODY BEGIN (Sections 5-8)
% =====================================================================
\Needspace{16\baselineskip}
\section{\texorpdfstring{Using the Lens: Dissect or Self-Check\added}{Using the Lens: Dissect or Self-Check}}

The lens has two modes. Choose one before you start.

\begin{panel}[breakable=false]
\textbf{Dissect} a paper you are reading: record what it does, and where, for fields \lensid{1}--\lensid{10}.
A dissection describes the paper; it does not grade it.\par\smallskip
\textbf{Self-check} a draft you are writing: give every field, \lensid{1}--\lensid{12}, a status and one fix
move.\par\smallskip
\textbf{Status key, self-check only:}\enspace \stS{} does it\quad \stP{} partly\quad \stM{} not at all\quad
\stN{} n/a\par\smallskip
\textbf{Locations:} section, paragraph and the sentence's first words, e.g.\ \loc{Sec.~2 P3 ``Existing EV and
PT\ldots''}; figures and tables by number, e.g.\ \loc{Fig.~1}.
\end{panel}

\begin{lensfields}
  \item[R1] \textbf{Choose the mode first.} Dissect a paper you are reading; self-check a draft you are writing.
  \item[R2] \textbf{Every finding cites a location,} written so the reader can find it: section, paragraph and the
        sentence's first words. Where sentences are numbered, as in Section~6, sentence numbers do.
  \item[R3] \textbf{A dissection describes; it never grades.} ``Not present'' and ``stated only in Sec.~5'' are facts
        and may be reported; statuses and verdicts may not.
  \item[R4] \textbf{A self-check gives each field one status and at most one fix move,} names the method behind
        the fix (\methodid{1}--\methodid{10}), and leaves numbers, results and technical claims untouched.
  \item[R5] \textbf{Quotes stay short.} At most one quoted sentence per finding, trimmed to about 25 words.
  \item[R6] \textbf{Read the whole paper.}\added{} Start with the introduction, then read every section, and record
        each finding wherever it is stated; its location shows which section that is.
\end{lensfields}

% =====================================================================
\section{\texorpdfstring{Worked Example\added}{Worked Example}}

The same short introduction is dissected below as if someone else wrote it, then self-checked as if it were your
draft. It is a teaching text written for this handout: the paper and its Figure~1 are invented, while the
background in P1--P2 is standard in the LES literature. It is an introduction only, so every location below is
in the introduction.

\begin{panel}[breakable=false]
\pn{1}\sn{1}Most LES practitioners close the subgrid stress with an eddy viscosity, and in doing so they assume
that energy flows only from the resolved scales to the subgrid scales.
\sn{2}A priori tests on filtered DNS data have long shown that this assumption is incomplete: locally, energy
also flows back to the resolved scales.
\sn{3}Because a positive eddy viscosity is purely dissipative, the model cannot represent this backscatter, and
where backscatter is strong the closure drains energy from scales that are in fact gaining it.
\sn{4}Figure~1 makes the mismatch visible: the modeled flux is forward everywhere, while the filtered-DNS flux
changes sign across the domain.\par\medskip
\pn{2}\sn{1}Two lines of work have addressed backscatter, and each leaves the same question open.
\sn{2}Stochastic models add a random forcing whose level is set to match an average backscatter rate; matching
an average, however, says little about where reverse transfer occurs.
\sn{3}Dynamic procedures can return a negative eddy viscosity, but the resulting instability is commonly removed
by clipping or averaging, which discards the very events at issue.
\sn{4}Neither approach tells a modeler where their closure is wrong.\par\medskip
\pn{3}\sn{1}This paper gives modelers a tool for answering that question: a local diagnostic that measures the
sign and size of interscale energy transfer at each point of a filtered field.
\sn{2}Applied to filtered DNS of decaying isotropic turbulence, it shows where an eddy-viscosity closure removes
energy that the flow is gaining, so modelers can test their closure where it is most likely to fail.
\sn{3}Section~2 defines the diagnostic; Section~3 applies it.
\end{panel}

\subsection{Dissect (fields 1--10)}

{\footnotesize
\renewcommand{\arraystretch}{1.1}
\begin{tabularx}{\linewidth}{@{}>{\raggedright\arraybackslash\sffamily\bfseries}p{2.5cm}>{\raggedright\arraybackslash}X>{\raggedright\arraybackslash}p{2.3cm}@{}}
\toprule
\normalfont\sffamily\bfseries Field & \sffamily\bfseries What the paper does & \sffamily\bfseries Where \\
\midrule
1- Readers      & Research community: LES modelers who close the subgrid stress with an eddy viscosity, plus
                  stochastic modelers and users of dynamic procedures. Wider readers: not named. Venue implied:
                  an LES or turbulence-modeling journal & \loc{P1 S1; P2 S2--S3} \\
2- Value words  & Instability, e.g.\ ``this assumption is \textbf{incomplete}''; ``the model \textbf{cannot} represent
                  this backscatter''; ``matching an average, \textbf{however}, says little''; ``\textbf{Neither}
                  approach tells a modeler where their closure is \textbf{wrong}.'' Community, e.g.\ ``\textbf{Most}
                  LES practitioners''; ``have \textbf{long shown}''; ``is \textbf{commonly} removed.'' Frame:
                  ``[Evidence] has long shown that this assumption is incomplete''
                & \loc{P1 S2; P1 S3; P2 S2; P2 S4; P1 S1; P1 S2; P2 S3} \\
3- Value offered & ``so modelers can test their closure where it is most likely to fail.'' Label (interpretation):
                  a diagnostic method & \loc{P3 S2} \\
4- Problem      & Current state: eddy-viscosity closures assume energy flows only to the subgrid scales. Where it
                  breaks: a dissipative closure cannot represent backscatter (Fig.~1); stochastic models match only
                  an average rate; dynamic procedures clip the negative events. Open: where is a modeler's closure
                  wrong? Explicit & \loc{P1 S1; P1 S2--S4; P2 S2--S3; P2 S4} \\
5- Stance       & Direct toward the dominant practice (``incomplete''); credits both rival lines of work before
                  naming their shared limit & \loc{P1 S2; P2 S1--S4} \\
6- Stakes       & Cost: the closure drains energy from scales that are gaining it, stated as physics. Benefit:
                  closures can be tested where they are most likely to fail & \loc{P1 S3; P3 S2} \\
7- Doubts answered & Not present in the introduction & -- \\
8- Contribution & A local diagnostic of the sign and size of interscale transfer, named as a tool & \loc{P3 S1} \\
9- Structure    & Problem (P1), prior work that sharpens it (P2), contribution (P3); explanation deferred to
                  Section~2; error model: both approaches share a flaw & \loc{P1--P3; P3 S3; P2 S4} \\
10- Pitfalls    & None of the four checks found. Background opening: the first instability word,
                  ``incomplete,'' is in P1. No aim or interest sentence; no novelty claim; results without
                  consequence does not apply to an introduction & \loc{P1 S2} \\
\bottomrule
\end{tabularx}
\renewcommand{\arraystretch}{1}
}

\clearpage
\subsection{Self-check (fields 1--12)}

{\footnotesize
\renewcommand{\arraystretch}{1.1}
\begin{tabularx}{\linewidth}{@{}>{\raggedright\arraybackslash\sffamily\bfseries}p{2.5cm}l>{\raggedright\arraybackslash}X@{}}
\toprule
\normalfont\sffamily\bfseries Field & \sffamily\bfseries Status & \sffamily\bfseries Evidence and fix move \\
\midrule
1- Readers      & \stS & Research community named in the first words; rival communities named in P2
                         (\loc{P1 S1; P2 S2--S3}) \\
2- Value words  & \stS & Instability words in every paragraph of the problem (\loc{P1 S2--S3; P2 S2--S4}) \\
3- Value offered & \stS & Offered as something modelers can do (\loc{P3 S2}; M1) \\
4- Problem      & \stS & All three parts stated: current state, where it breaks, what remains open
                         (\loc{P1 S1--S4; P2 S2--S4}) \\
5- Stance       & \stS & Credits both lines of work before naming their shared limit (\loc{P2 S1--S4}; M9) \\
6- Stakes       & \stP & The cost is stated as physics, not as a cost to these readers. Say what misplaced transfer
                         does to the predictions they care about (M2, M10) \\
7- Doubts answered & \stM & No doubt is anticipated. Name the doubt modelers will raise, such as whether an a priori
                         diagnostic says anything about a posteriori runs, and say where the paper answers it (M4) \\
8- Contribution & \stS & Named as a tool these readers can use (\loc{P3 S1}; M1) \\
9- Structure    & \stS & Problem, then prior work that sharpens it, then contribution; definitions deferred to
                         Section~2 (M8, M7) \\
10- Pitfalls    & \stS & None of the four checks found; the results check does not apply to an introduction \\
11- Value count & \stS & Nine distinct value words in P1--P2 \\
12- Opening     & \stP & \loc{P1 S1} states the readers' practice, but the question the paper answers (where is my
                         closure wrong?) first appears in \loc{P2 S4}. Bring a version of it into P1 (M8) \\
\bottomrule
\end{tabularx}
\renewcommand{\arraystretch}{1}
}

{\small\textbf{Priority fix:} 6- Stakes. Readers can see that the model is wrong, but not yet what it costs them.
Doubts answered (7-) is missing too, but readers who can't see the cost stop reading before their doubts matter.\par}

% =====================================================================
\clearpage
\section{\texorpdfstring{Worksheet\added}{Worksheet}}

\renewcommand{\arraystretch}{1.6}
\begin{tabularx}{\linewidth}{@{}>{\sffamily\small\bfseries}p{3.4cm}X@{}}
Text analyzed       & \hrulefill \\
Mode                & Dissect \enspace / \enspace Self-check \\
Readers (L1)        & \hrulefill \\
\end{tabularx}

\renewcommand{\arraystretch}{1.9}
{\small
\begin{tabularx}{\linewidth}{@{}|l|X|p{2.3cm}|p{1.5cm}|X|@{}}
\hline
\sffamily\bfseries ID & \sffamily\bfseries Finding (what the text does) & \sffamily\bfseries Location
  & \sffamily\bfseries Status & \sffamily\bfseries Fix move (self-check) \\ \hline
L1  & & & & \\ \hline
L2  & & & & \\ \hline
L3  & & & & \\ \hline
L4  & & & & \\ \hline
L5  & & & & \\ \hline
L6  & & & & \\ \hline
L7  & & & & \\ \hline
L8  & & & & \\ \hline
L9  & & & & \\ \hline
L10 & & & & \\ \hline
L11 & & & & \\ \hline
L12 & & & & \\ \hline
\end{tabularx}
}
\renewcommand{\arraystretch}{1}

{\small Status is for self-check only; leave it blank when dissecting (R3). L11 and L12 apply to your own draft.}

\textbf{Reusable frames} (L2): \hrulefill\par
\vspace{4pt}\hrulefill\par
\textbf{Priority fix} (self-check): \hrulefill

% =====================================================================
\Needspace{27\baselineskip}
\section{\texorpdfstring{Report Format for Skills and Reviewers\added}{Report Format for Skills and Reviewers}}

Reports produced from this handout, by a person or by a skill, use this layout so they can be compared and
cited. A skill writes it in ordinary chat formatting: bold labels and a table, not a code block.

\begin{panel}[breakable=false]
\small
\begin{verbatim}
McEnerney-lens report - Dissect | Self-check
Text      : title, venue, sections read
Readers   : research community; wider readers
Labels    : "1- Readers" to "12- Opening", never the IDs; cells up to
            about 60 words, "4- Problem" about 40 for each of its parts

DISSECT    : table | Field | What the paper does | Where |, rows 1- to 10-,
             then tables "2- Value words in context" and "10- Pitfalls"
             "Not present" and "stated only in Sec. n" are facts; no statuses

SELF-CHECK : the same two tables, then | Field | Status | Evidence and fix move |
             one row per field, 1- to 12-: STRONG|PARTIAL|MISSING|N/A,
             evidence with its location, fix (M<n>)
             Priority fix: the one field that most limits value, and why
             Untouched   : numbers, results and technical claims are not changed

Where     : section, paragraph and the sentence's first words (R2)
\end{verbatim}
\end{panel}
% <<< ADDITIONS-BODY END

% =====================================================================
\section{References}
\enlargethispage{1\baselineskip}% keeps the closing notes on the same page

{\small
\renewcommand{\section}[2]{}%  suppress the bibliography's own heading
\begin{thebibliography}{9}\setlength{\itemsep}{0pt}\raggedright

\bibitem{mcenerney2014}
L.~McEnerney, \emph{Leadership Lab: The Craft of Writing Effectively}, lecture, University of Chicago Writing
Program. Recording posted by UChicago Social Sciences on YouTube, 26 June 2014.
\url{https://youtu.be/vtIzMaLkCaM}

\bibitem{sewell1996}
W.~H. Sewell Jr., ``Historical events as transformations of structures: Inventing revolution at the Bastille,''
\emph{Theory and Society}, 25(6), 841--881, 1996. doi:10.1007/BF00159818. (The ``Ever since Herodotus''
opening: a problem built in the first paragraph, with its two reader communities named.)

% TODO(verify): end page of the Tutino article; only the volume,
% issue and first page (78(3), 367) were confirmed when this was compiled.
\bibitem{tutino1998}
J.~Tutino, ``The revolution in Mexican independence: Insurgency and the renegotiation of property, production,
and patriarchy in the Baj\'{\i}o, 1800--1855,'' \emph{Hispanic American Historical Review}, 78(3), 367ff.,
August 1998. (An opening that directly challenges an entrenched view; its first footnote shows readers
helping the author find his problem.)

\bibitem{said1978}
E.~W. Said, \emph{Orientalism}, Pantheon Books, 1978. (Its introduction is McEnerney's extended example of a
writer defining several communities of readers.)

\bibitem{lyotard1984}
J.-F. Lyotard, \emph{The Postmodern Condition: A Report on Knowledge} (1979), trans. G.~Bennington and
B.~Massumi, Manchester University Press, 1984, section~1. (The passage on knowledge becoming exterior to the
knower and taking the form of a commodity.)

\bibitem{wittgenstein1953}
L.~Wittgenstein, \emph{Philosophical Investigations}, Blackwell, 1953. (McEnerney's ``what something means is
what it does'': meaning as use.)

\bibitem{vrijenhoek1979}
R.~C. Vrijenhoek, ``Factors affecting clonal diversity and coexistence,'' \emph{American Zoologist}, 19,
787--797, 1979. (Background to the frozen niche-variation model named in the lecture's 1A/1B biology example;
the two example paragraphs are teaching texts.)

\bibitem{booth}
W.~C. Booth, G.~G. Colomb and J.~M. Williams, \emph{The Craft of Research}, University of Chicago Press (4th ed.,
2016, with J.~Bizup and W.~T. FitzGerald). (Further reading, not named in the lecture: Colomb and Williams
developed the course McEnerney teaches, and the book treats research problems at length.)

\end{thebibliography}

\textbf{Not identified in the recording:} the ``why people write essays'' passage McEnerney criticizes (he calls
it a caricature written for high-school students) \tsin{44:29--49:22}, and the economists' introduction he praises
on pages 16--17 of his printed handout \tsin{1:14:34--1:15:34}.

\textbf{See also:} Handout 2 (\emph{effective-writing}) applies the same reader-value idea to single paragraphs.
}

\end{document}
